\documentclass[12pt]{article}

% Packages
\usepackage[a4paper, margin=1in]{geometry}
\usepackage{setspace}
\usepackage{titlesec}
\usepackage{enumitem}
\usepackage{hyperref}

% Formatting
\onehalfspacing
\setlist[itemize]{noitemsep, topsep=0pt}
\setlist[enumerate]{noitemsep, topsep=0pt}

\titleformat{\section}{\large\bfseries}{\thesection.}{0.5em}{}
\titleformat{\subsection}{\normalsize\bfseries}{\thesubsection}{0.5em}{}
\titleformat{\subsubsection}{\normalsize\itshape}{\thesubsubsection}{0.5em}{}

\begin{document}

\begin{center}
    {\Large \textbf{CSE400 -- Lecture 3 Scribe}}\\[0.3cm]
    {\large Project Kickoff and Undergraduate Research Programme (UGRP)}
\end{center}

\vspace{0.5cm}

\section{Course Project Component}

\subsection{Definition}
The CSE400 course includes a \textbf{Project Kickoff} component carrying a weightage of \textbf{30\%} of the total course evaluation.

\subsection{Assumption}
Since the project carries significant weightage, it is intended to be a core learning and evaluation mechanism rather than an auxiliary activity.

\section{Project Team Formation}

\subsection{Requirement}
\textbf{Deadline:} 17 January 2026 (Saturday, End of Day)

\subsection{Reasoning}
\begin{itemize}
    \item The project is group-based.
    \item Group-based execution requires early coordination.
    \item A fixed team formation deadline ensures timely initiation of project milestones.
\end{itemize}

\section{Project Execution Guidelines}

\subsection{Weightage}
Total project weightage: \textbf{30\%}

\subsection{Structure}
\begin{itemize}
    \item The project consists of six major milestones (M1--M6).
    \item One submission per group is required for each milestone.
\end{itemize}

\subsection{Assumption}
The milestones are sequential and cumulative, meaning each stage builds upon the previous one.

\section{Milestone-Wise Project Execution}

\subsection{M1: Kickstart}

\subsubsection{Components}
\begin{itemize}
    \item Team Formation
    \item Area Identification
    \item Background
    \item Motivation
    \item Problem Formulation
\end{itemize}

\subsubsection{Step-by-Step Reasoning}
\begin{itemize}
    \item A project must begin with a clearly identified domain.
    \item Understanding the background provides context.
    \item Motivation justifies the importance of the problem.
    \item Based on background and motivation, the problem is formally formulated.
\end{itemize}

\subsection{M2: Mathematical Modeling}

\subsubsection{Definition}
Mathematical modeling of the selected problem using:
\begin{itemize}
    \item Random Variables (RV)
    \item PMF / PDF
    \item CDF
    \item Multivariate Random Variables
    \item Joint PMF / PDF / CDF
\end{itemize}

\subsubsection{Step-by-Step Reasoning}
\begin{itemize}
    \item Real-world problems contain uncertainty.
    \item Uncertainty is modeled using random variables.
    \item PMF/PDF describes probabilistic behavior.
    \item CDF captures cumulative probability.
    \item Joint distributions are required for multiple variables.
\end{itemize}

\subsection{M3: Coding}

\subsubsection{Components}
\begin{itemize}
    \item Simulation
    \item Computation
\end{itemize}

\subsubsection{Reasoning}
\begin{itemize}
    \item Mathematical models must be validated.
    \item Simulation reproduces model behavior.
    \item Computation provides numerical results.
\end{itemize}

\subsection{M4: Inference}

\subsubsection{Components}
\begin{itemize}
    \item Choose a randomized algorithm
    \item Understand the algorithm
    \item Code the algorithm
\end{itemize}

\subsubsection{Assumption}
Randomized algorithms are suitable for inference in the modeled problem domain.

\subsubsection{Reasoning}
\begin{itemize}
    \item Algorithm selection precedes implementation.
    \item Understanding ensures correctness.
    \item Coding enables evaluation.
\end{itemize}

\subsection{M5: Randomized Algorithms}

\subsubsection{Components}
\begin{itemize}
    \item Apply the randomized algorithm
    \item Generate results
    \item Compare with deterministic algorithms
\end{itemize}

\subsubsection{Step-by-Step Reasoning}
\begin{itemize}
    \item Apply the algorithm to the problem.
    \item Obtain execution results.
    \item Compare with deterministic methods.
    \item Draw new inferences.
\end{itemize}

\subsection{M6: Bounds and Analysis}

\subsubsection{Components}
\begin{itemize}
    \item Derive bounds
    \item Perform analysis
    \item Compile final deliverables
\end{itemize}

\subsubsection{Reasoning}
\begin{itemize}
    \item Theoretical bounds validate performance.
    \item Analysis supports empirical results.
    \item All outputs are compiled for final submission.
\end{itemize}

\section{Project Deliverables and Assessment}

\subsection{Deliverables}
\begin{itemize}
    \item Codes
    \item Reports
    \item Videos
    \item Other milestone-specific artifacts
\end{itemize}

\subsection{Assessment}
\begin{itemize}
    \item Team assessment before midsemester
    \item Team assessment after midsemester
    \item Final evaluation:
    \begin{itemize}
        \item Project Viva
        \item Final Submission
    \end{itemize}
\end{itemize}

\section{Submission \#1: Concept Evolution Maps}

\subsection{Definition}
Concept evolution maps visualize how project ideas evolve from initial concepts to final problem formulation.

\subsection{Tools}
\begin{itemize}
    \item Miro
    \item draw.io (diagrams.net)
\end{itemize}

\section{Submission \#2: Scribe -- Learning Reflection Logs}

\subsection{Types}
\begin{itemize}
    \item Lecture Scribe
    \item Project Scribe
\end{itemize}

\subsection{Lecture Scribe}
\begin{itemize}
    \item Assigned to two groups per lecture
    \item Reflects lecture content
    \item Minimum length: 8--10 pages
\end{itemize}

\subsection{Project Scribe Components}
\begin{itemize}
    \item Decision logs (Why X over Y?)
    \item Constraints and alternatives
    \item Final decisions
    \item Evidence and justification
    \item Trade-off matrices (Cost vs Performance vs Risk)
\end{itemize}

\section{Submission \#3: Multimodal Artifacts}

\subsection{Definition}
Artifacts may include video, audio, or visual materials.

\subsection{Requirements}
\begin{itemize}
    \item One video per milestone
    \item Duration: 10--15 minutes
    \item Focus on content, not editing
\end{itemize}

\subsection{Types}
\begin{itemize}
    \item Think-Aloud Videos
    \item One-Minute Insight Videos
    \item Project Demo
\end{itemize}

\section{Introduction to UGRP-8 (2026--2027)}

UGRP is introduced as a structured undergraduate research initiative.

\section{Motivation for UGRP}

\subsection{T-Shaped Engineer Concept}
\textbf{Vertical bar:} Depth in a single technical discipline\\
\textbf{Horizontal bar:} Ability to apply knowledge across disciplines and collaborate

\subsection{Reasoning}
UGRP aims to develop both technical depth and interdisciplinary breadth.

\section{Philosophy of UGRP}

\subsection{Core Principles}
\begin{itemize}
    \item Multidisciplinary
    \item Experiential Learning
    \item Research Driven
\end{itemize}

\subsection{4D Model}
\begin{enumerate}
    \item Discover
    \item Design
    \item Develop
    \item Deliver
\end{enumerate}

\section{Technical Breadth Under UGRP}

Includes:
\begin{itemize}
    \item CS and CSE
    \item Data Science and Applied AI
    \item Modern Computer System Design
    \item Networks
    \item IoT / IoBNT / IoV
\end{itemize}

\section{Industry and Research Activities}

\subsection{Example}
5G-enabled Intelligent Transportation Systems Testbed in Gujarat

\subsection{Activities}
\begin{itemize}
    \item V2X communication
    \item Sensor data collection
    \item ITS dashboard development
    \item Deep learning and deep-RL solutions
    \item mmWave beam prediction
    \item C-V2X and Wi-Fi coexistence
\end{itemize}

\section{Collaborations and Outcomes}

\subsection{Evidence Presented}
\begin{itemize}
    \item International academic collaborations
    \item IEEE conference and journal publications
    \item Awards and recognitions
    \item Alumni placements in industry and academia
\end{itemize}

\section{Conclusion of Lecture 3}

Lecture 3 establishes:
\begin{itemize}
    \item The project execution framework of CSE400
    \item A milestone-driven research methodology
    \item The UGRP philosophy and its academic and industrial outcomes
\end{itemize}

The lecture emphasizes process, structure, and research orientation, forming the foundation for subsequent technical work in the course.

\end{document}
